% ID: SEL-D0072-OA-E03
% ESTADO: cotejada. Segunda pasada visual y borrador compilado revisados.
% Fuente: Selectividad/sel_2009_matematicas_aplicadas/mat apli CC SS06.pdf
% SHA-256: b371bdd6ba8b176301371f25945c7fb030327ff4501a6a02aa6f109a3a4f53e1
% Página: 1; opción A; ejercicio original 3.

\noindent\underline{Parte I}

Se consideran dos sucesos $A$ y $B$, asociados a un espacio muestral, tales que
\[
P(A\cup B)=1,\quad P(A\cap B)=0.3\quad\text{y}\quad P(A/B)=0.6.
\]
\begin{enumerate}[label=\alph*)]
\item \textbf{(1.5 puntos)} Halle las probabilidades de los sucesos $A$ y $B$.
\item \textbf{(0.5 puntos)} Determine si el suceso $B$ es independiente del suceso $A$.
\end{enumerate}
\par\noindent\begin{minipage}{\linewidth}
\noindent\underline{Parte II}

El gasto que hacen las familias españolas en regalos de Navidad sigue una ley Normal de media desconocida y desviación típica 84 euros. Para estimar esta media se seleccionó una muestra aleatoria y se obtuvo el intervalo de confianza $(509.41,539.79)$, con un nivel de confianza del $97\%$.
\begin{enumerate}[label=\alph*)]
\item \textbf{(0.5 puntos)} \textquestiondown{}Cuál ha sido la media de la muestra escogida?
\item \textbf{(1.5 puntos)} \textquestiondown{}Qué tamaño tenía la muestra?
\end{enumerate}
\end{minipage}\par
